\begin{gradedsolution}
\begin{enumerate}
    \item Let $\fun \setin \funspop$, and let $\res, \EM{{\colR s}} \setin \ressp$ with $\res \posleq_{\ressp} \EM{{\colR s}}$.
    Since $\fun \posleq_{\funspop} \fun$ (reflexivity), we have $\tup{\fun, \res} \posleq_{\funspop \Ptimes \ressp} \tup{\fun, \EM{{\colR s}}}$ in the product order.
    Since $\adp$ is monotone, it follows that $\adp(\tup{\fun, \res}) \leq \adp(\tup{\fun, \EM{{\colR s}}})$.
    This shows that the map $\ressp \mto \Bool$, $\res \mapsto \adp(\tup{\fun, \res})$ is monotone.
    The subset $\widehat{\adp}(\fun) \setsubseteq \ressp$ consists precisely of those $\res$ which this monotone map sends to $\true$, so by the first part of \cref{ex:UpperSetsAsMonotoneMaps}, it is an upper set of $\ressp$.

    \item To show that $\widehat{\adp}$ is a monotone map, we need to show that whenever $\fun, \EM{{\colF g}} \setin \funspop$ with $\fun \posleq_{\funspop} \EM{{\colF g}}$, then $\widehat{\adp}(\fun) \setsubseteq \widehat{\adp}(\EM{{\colF g}})$.
    Suppose $\fun, \EM{{\colF g}} \setin \funspop$ with $\fun \posleq_{\funspop} \EM{{\colF g}}$, and let $\res \setin \widehat{\adp}(\fun)$ be arbitrary.
    Since $\tup{\fun, \res} \posleq_{\funspop \Ptimes \ressp} \tup{\EM{{\colF g}}, \res}$ and $\adp$ is monotone, it follows that $\adp(\tup{\fun, \res}) \leq \adp(\tup{\EM{{\colF g}}, \res})$.
    Since $\adp(\tup{\fun, \res}) = \true$ and $\true$ is the greatest element of $\Bool$, it must follow that $\adp(\tup{\EM{{\colF g}}, \res}) = \true$.
    Hence $\res \setin \widehat{\adp}(\EM{{\colF g}})$, and since $\res \setin \widehat{\adp}(\fun)$ was chosen arbitrarily, we have shown that $\widehat{\adp}(\fun) \setsubseteq \widehat{\adp}(\EM{{\colF g}})$.
    This shows that $\widehat{\adp}$ is a monotone map.

    \item Let $\fun, \EM{{\colF g}} \setin \funsp$ with $\fun \posleq_{\funsp} \EM{{\colF g}}$.
    By definition of the opposite poset, this means $\EM{{\colF g}} \posleq_{\funspop} \fun$, and so, by the previous part, $\widehat{\adp}(\EM{{\colF g}}) \setsubseteq \widehat{\adp}(\fun)$.

    Interpretation: any resources that suffice to provide the functionality $\EM{{\colF g}}$ also suffice to provide the less demanding functionality $\fun$.
    Equivalently: demanding more functionality can only shrink, never enlarge, the set of feasible resources.
\end{enumerate}
\end{gradedsolution}
